\subsection{Command}
The \texttt{Command} class represents the fundamental unit of execution
in the EOIW framework. Each command encapsulates the metadata required for
tracking, auditing, and ensuring idempotency during processing \shortcite{hopcroft1983data}.

\begin{algorithm}[H]
\caption{Command (\emph{Class})}
\begin{algorithmic}[1]
\Procedure{Command}{Command\textlangle T\textrangle}
    \State \textbf{Fields:}
    \Statex \Comment{Metadata fields for tracking and auditing}
    \State \quad id : UUID
    \State \quad idempotencyKey : String
    \Statex \Comment{Status and timestamps}
    \State \quad status : String
    \State \quad createdAt : DateTime
    \State \quad updatedAt : DateTime
    \Statex \Comment{Tenant and user context}
    \State \quad tenantId : String
    \State \quad username : String
    \Statex \Comment{Generic payload}
    \State \quad payload : T
\EndProcedure
\end{algorithmic}
\end{algorithm}

\noindent
This class serves as a base abstraction for all executable operations \shortcite{servlets2002java}. Each
subclass specializes the \texttt{payload} field to represent domain-specific data.